\documentclass[10pt,letterpaper]{article}
\usepackage[T1]{fontenc}
\usepackage[utf8]{inputenc}
\usepackage[spanish,es-tabla]{babel}
\usepackage{csquotes}
\usepackage{mathpazo}
\usepackage[scaled=0.9]{helvet}
\usepackage{courier}
\usepackage{calc}
\usepackage[top=2cm, bottom=2cm, left=2cm, right=2cm]{geometry}
\usepackage{xcolor}
\usepackage{booktabs}
\usepackage{longtable}
\usepackage{array}
\usepackage{ragged2e}
\usepackage{xurl}
\usepackage[strings]{underscore}
\usepackage{fancyhdr}
\usepackage{sectsty}
\usepackage{tocloft}
\usepackage[colorlinks=true,linkcolor=blue,urlcolor=blue,citecolor=blue,breaklinks=true]{hyperref}
\usepackage{enumitem}

\definecolor{uncpblue}{HTML}{1A237E}
\allsectionsfont{\color{uncpblue}}

\renewcommand{\cftsecleader}{\cftdotfill{\cftsecdotsep}}
\cftpagenumbersoff{section}
\cftpagenumbersoff{subsection}
\cftpagenumbersoff{subsubsection}
\renewcommand{\cftsecfont}{\normalfont}
\renewcommand{\cftsubsecfont}{\normalfont}
\renewcommand{\cftsubsubsecfont}{\normalfont}
\setlength{\parindent}{0pt}
\setlength{\parskip}{4pt}
\setlength{\tabcolsep}{3pt}
\renewcommand{\arraystretch}{0.85}

\pagestyle{fancy}
\fancyhf{}
\fancyhead[L]{\small\textcolor{uncpblue}{\textbf{SGSI-UNCP}}}
\fancyhead[R]{\small\textcolor{uncpblue}{P-SGSI-08}}
\fancyfoot[C]{\thepage}
\renewcommand{\headrulewidth}{0.4pt}
\renewcommand{\footrulewidth}{0.4pt}
\setlength{\headheight}{14pt}

\setlist{nosep,leftmargin=*,after=\vspace{2pt}}
\setlength{\emergencystretch}{2em}

\newcommand{\makecover}{
\begin{titlepage}
\centering
\vspace*{2cm}
{\Huge\bfseries\color{uncpblue} SGSI-UNCP\par}
\vspace{0.7cm}
{\LARGE Plan Anual de Capacitación y Concientización en Seguridad de la
Información\par}
\vspace{0.5cm}
{\large Documento Base del Sistema de Gestión\\de Seguridad de la Información\par}
\vspace{1.5cm}
{\small Universidad Nacional del Centro del Perú\par}
\vfill
\end{titlepage}
}

\begin{document}

\makecover
\setcounter{tocdepth}{2}
\RaggedRight
\tableofcontents
\newpage

\RaggedRight
\section{Plan Anual de Capacitación y Concientización en Seguridad de la
Información}\label{plan-anual-de-capacitaciuxf3n-y-concientizaciuxf3n-en-seguridad-de-la-informaciuxf3n}

\textbf{Organización:} Universidad Nacional del Centro del Perú (UNCP)
\textbf{Año:} 2026 \textbf{Responsable:} Oficial de Seguridad / Unidad
de RRHH \textbf{Presupuesto Estimado:} S/ 45,000 \textbf{Alineamiento:}
OGTD5 (Fortalecer la gestión de competencias digitales), OGTD6
(Desarrollar competencias digitales), POL-SGSI-02 (Uso Aceptable)

\begin{center}\rule{0.5\linewidth}{0.5pt}\end{center}

\subsection{Objetivos}\label{objetivos}

\subsubsection{1.1 Generales}\label{generales}

\begin{itemize}
\tightlist
\item
  Elevar el nivel de cultura de seguridad de la información en toda la
  comunidad universitaria (personal administrativo, docente, autoridades
  y estudiantes).
\item
  Reducir la probabilidad y el impacto de incidentes de seguridad
  causados por error humano (phishing, pérdida de credenciales, uso
  inadecuado de recursos).
\item
  Desarrollar las competencias técnicas del personal de la OTI para
  operar y mantener los controles de seguridad implementados en el marco
  del SGSI.
\end{itemize}

\subsubsection{1.2 Específicos (Metas
2026)}\label{especuxedficos-metas-2026}

{\setlength{\RaggedRightRightskip}{0pt plus 5em}\small\sloppy \begin{longtable}[]{@{}
  >{\hspace{0pt}\RaggedRight\arraybackslash}p{(\linewidth - 6\tabcolsep) * \real{0.2500}}
  >{\hspace{0pt}\RaggedRight\arraybackslash}p{(\linewidth - 6\tabcolsep) * \real{0.2500}}
  >{\hspace{0pt}\RaggedRight\arraybackslash}p{(\linewidth - 6\tabcolsep) * \real{0.2500}}
  >{\hspace{0pt}\RaggedRight\arraybackslash}p{(\linewidth - 6\tabcolsep) * \real{0.2500}}@{}}
\toprule\noalign{}
\begin{minipage}[b]{\linewidth}\raggedright
Meta
\end{minipage} & \begin{minipage}[b]{\linewidth}\raggedright
Indicador
\end{minipage} & \begin{minipage}[b]{\linewidth}\raggedright
Línea Base
\end{minipage} & \begin{minipage}[b]{\linewidth}\raggedright
Meta
\end{minipage} \\
\midrule\noalign{}
\endhead
\bottomrule\noalign{}
\endlastfoot
Cobertura de capacitación obligatoria & \% de trabajadores que completan
el curso anual de fundamentos de ciberseguridad & 0\% (nuevo programa) &
100\% \\
Reducción de vulnerabilidad al phishing & \% de usuarios que hacen clic
en enlaces de simulacros de phishing & Esperado \textasciitilde30\%
inicial & \textless{} 10\% al final del año \\
Personal técnico certificado & Número de certificaciones de seguridad
obtenidas por el equipo OTI & 0 & 3 certificaciones \\
Satisfacción de la capacitación & Calificación promedio de los cursos
(escala 1-5) & N/A & \textgreater= 4.0 \\
\end{longtable}\normalsize}

\begin{center}\rule{0.5\linewidth}{0.5pt}\end{center}

\subsection{Programa de Capacitación
Técnica}\label{programa-de-capacitaciuxf3n-tuxe9cnica}

\emph{Dirigido al personal de la OTI, administradores de sistemas y
responsables de seguridad.}

{\setlength{\RaggedRightRightskip}{0pt plus 5em}\small\sloppy \begin{longtable}[]{@{}
  >{\hspace{0pt}\RaggedRight\arraybackslash}p{(\linewidth - 10\tabcolsep) * \real{0.1667}}
  >{\hspace{0pt}\RaggedRight\arraybackslash}p{(\linewidth - 10\tabcolsep) * \real{0.1667}}
  >{\hspace{0pt}\RaggedRight\arraybackslash}p{(\linewidth - 10\tabcolsep) * \real{0.1667}}
  >{\hspace{0pt}\RaggedRight\arraybackslash}p{(\linewidth - 10\tabcolsep) * \real{0.1667}}
  >{\hspace{0pt}\RaggedRight\arraybackslash}p{(\linewidth - 10\tabcolsep) * \real{0.1667}}
  >{\hspace{0pt}\RaggedRight\arraybackslash}p{(\linewidth - 10\tabcolsep) * \real{0.1667}}@{}}
\toprule\noalign{}
\begin{minipage}[b]{\linewidth}\raggedright
Tema
\end{minipage} & \begin{minipage}[b]{\linewidth}\raggedright
Modalidad
\end{minipage} & \begin{minipage}[b]{\linewidth}\raggedright
Duración (horas)
\end{minipage} & \begin{minipage}[b]{\linewidth}\raggedright
Mes
\end{minipage} & \begin{minipage}[b]{\linewidth}\raggedright
Proveedor Sugerido
\end{minipage} & \begin{minipage}[b]{\linewidth}\raggedright
Costo Estimado (S/)
\end{minipage} \\
\midrule\noalign{}
\endhead
\bottomrule\noalign{}
\endlastfoot
Implementación y Auditoría de ISO 27001:2022 & Virtual / Presencial & 24
hrs (3 días) & Marzo & Consultora especializada & 12,000 \\
Administración Segura en Huawei Cloud & Virtual (oficial Huawei) & 20
hrs (5 sesiones) & Abril & Huawei Academy / Partner & 8,000 \\
Gestión de Incidentes y Respuesta a Incidentes (CSIRT) & Presencial & 24
hrs (3 días) & Junio & Entrenador local o virtual & 10,000 \\
Desarrollo Seguro de Software (DevSecOps / SAST / DAST) & Virtual & 16
hrs (4 sesiones) & Agosto & Plataforma de entrenamiento técnico &
6,000 \\
Administración de SIEM (Wazuh / Sentinel) y Threat Hunting & Presencial
& 16 hrs (2 días) & Octubre & Consultora especializada & 9,000 \\
\textbf{Total Presupuesto Técnico} & & \textbf{100 hrs} & & &
\textbf{45,000} \\
\end{longtable}\normalsize}

\textbf{Certificaciones objetivo para el personal OTI (2026-2027):} -
ISO/IEC 27001 Lead Implementer o Lead Auditor (1 persona). - Huawei
Cloud Security Certification o equivalent (2 personas). - Certified
Incident Handler (EC-Council / SANS) o equivalent (1 persona).

\begin{center}\rule{0.5\linewidth}{0.5pt}\end{center}

\subsection{Programa de Concientización
General}\label{programa-de-concientizaciuxf3n-general}

\emph{Dirigido a todo el personal administrativo, docentes y
autoridades.}

\subsubsection{3.1 Actividades
Permanentes}\label{actividades-permanentes}

{\setlength{\RaggedRightRightskip}{0pt plus 5em}\small\sloppy \begin{longtable}[]{@{}
  >{\hspace{0pt}\RaggedRight\arraybackslash}p{(\linewidth - 8\tabcolsep) * \real{0.2000}}
  >{\hspace{0pt}\RaggedRight\arraybackslash}p{(\linewidth - 8\tabcolsep) * \real{0.2000}}
  >{\hspace{0pt}\RaggedRight\arraybackslash}p{(\linewidth - 8\tabcolsep) * \real{0.2000}}
  >{\hspace{0pt}\RaggedRight\arraybackslash}p{(\linewidth - 8\tabcolsep) * \real{0.2000}}
  >{\hspace{0pt}\RaggedRight\arraybackslash}p{(\linewidth - 8\tabcolsep) * \real{0.2000}}@{}}
\toprule\noalign{}
\begin{minipage}[b]{\linewidth}\raggedright
Actividad
\end{minipage} & \begin{minipage}[b]{\linewidth}\raggedright
Descripción
\end{minipage} & \begin{minipage}[b]{\linewidth}\raggedright
Canal
\end{minipage} & \begin{minipage}[b]{\linewidth}\raggedright
Frecuencia
\end{minipage} & \begin{minipage}[b]{\linewidth}\raggedright
Responsable
\end{minipage} \\
\midrule\noalign{}
\endhead
\bottomrule\noalign{}
\endlastfoot
Tips de Seguridad & Mensajes cortos con recomendaciones prácticas
(contraseñas seguras, phishing, cifrado, bloqueo de pantalla) & Correo
institucional & Mensual & Oficial de Seguridad \\
Video-Cápsulas ``Protege tus Datos'' & Videos de 2-3 minutos sobre temas
específicos de seguridad & Redes sociales UNCP + Moodle & Trimestral &
Oficial de Seguridad + Imagen Institucional \\
Simulacro de Phishing & Envío controlado de correos simulados para medir
y concientizar sobre identificación de phishing & Correo institucional &
Semestral (abril y octubre) & OTI (Seguridad) \\
Curso Virtual: Fundamentos de Ciberseguridad & Curso obligatorio con
evaluación final en el Campus Virtual & Moodle & Anual (permanente, con
recordatorio en marzo) & OTI + RRHH \\
Semana de la Ciberseguridad UNCP & Evento anual con charlas, talleres y
actividades interactivas & Presencial + Streaming & Anual (agosto) &
Comité organizador \\
\end{longtable}\normalsize}

\subsubsection{3.2 Curso Obligatorio: ``Fundamentos de Ciberseguridad
para el Personal
UNCP''}\label{curso-obligatorio-fundamentos-de-ciberseguridad-para-el-personal-uncp}

{\setlength{\RaggedRightRightskip}{0pt plus 5em}\small\sloppy \begin{longtable}[]{@{}
  >{\hspace{0pt}\RaggedRight\arraybackslash}p{(\linewidth - 4\tabcolsep) * \real{0.3333}}
  >{\hspace{0pt}\RaggedRight\arraybackslash}p{(\linewidth - 4\tabcolsep) * \real{0.3333}}
  >{\hspace{0pt}\RaggedRight\arraybackslash}p{(\linewidth - 4\tabcolsep) * \real{0.3333}}@{}}
\toprule\noalign{}
\begin{minipage}[b]{\linewidth}\raggedright
Módulo
\end{minipage} & \begin{minipage}[b]{\linewidth}\raggedright
Contenido
\end{minipage} & \begin{minipage}[b]{\linewidth}\raggedright
Duración Estimada
\end{minipage} \\
\midrule\noalign{}
\endhead
\bottomrule\noalign{}
\endlastfoot
Introducción a la Seguridad de la Información & Conceptos CID, activos
de información, SGSI UNCP & 30 min \\
Gestión de Contraseñas y MFA & Creación de contraseñas seguras, uso de
MFA, gestores de contraseñas & 20 min \\
Phishing e Ingeniería Social & Cómo identificar correos y mensajes
sospechosos, qué hacer ante un intento & 25 min \\
Uso Seguro de Correo y Dispositivos Móviles & Cifrado, VPN, bloqueo de
pantalla, POL-SGSI-05 aplicable al personal & 20 min \\
Protección de Datos Personales (Ley 29733) & Obligaciones del personal,
datos sensibles, consentimiento & 20 min \\
Reporte de Incidentes & Cuándo y cómo reportar, canales de comunicación
(incidentes-seguridad@uncp.edu.pe) & 15 min \\
\textbf{Evaluación Final} & 10 preguntas (mínimo 14/20 para aprobar) &
15 min \\
\textbf{Total} & & \textbf{2 horas 25 min} \\
\end{longtable}\normalsize}

\subsubsection{3.3 Campañas Especiales}\label{campauxf1as-especiales}

{\setlength{\RaggedRightRightskip}{0pt plus 5em}\small\sloppy \begin{longtable}[]{@{}
  >{\hspace{0pt}\RaggedRight\arraybackslash}p{(\linewidth - 4\tabcolsep) * \real{0.3333}}
  >{\hspace{0pt}\RaggedRight\arraybackslash}p{(\linewidth - 4\tabcolsep) * \real{0.3333}}
  >{\hspace{0pt}\RaggedRight\arraybackslash}p{(\linewidth - 4\tabcolsep) * \real{0.3333}}@{}}
\toprule\noalign{}
\begin{minipage}[b]{\linewidth}\raggedright
Campaña
\end{minipage} & \begin{minipage}[b]{\linewidth}\raggedright
Periodo
\end{minipage} & \begin{minipage}[b]{\linewidth}\raggedright
Tema Enfocado
\end{minipage} \\
\midrule\noalign{}
\endhead
\bottomrule\noalign{}
\endlastfoot
``Empieza Seguro'' & Marzo (inicio de clases) & Contraseñas, MFA,
bloqueo de pantalla \\
``Semana de la Ciberseguridad'' & Agosto & Charlas magistrales, talleres
prácticos, concurso de seguridad \\
``Cierre Seguro'' & Diciembre & Respaldo de información, cierre de
sesiones, seguridad en periodo vacacional \\
\end{longtable}\normalsize}

\begin{center}\rule{0.5\linewidth}{0.5pt}\end{center}

\subsection{Estrategia para Estudiantes (Alcance
Informativo)}\label{estrategia-para-estudiantes-alcance-informativo}

Dado que la UNCP no ejerce control técnico sobre los dispositivos
personales de los estudiantes, la estrategia se centra en la
\textbf{concientización voluntaria}:

{\setlength{\RaggedRightRightskip}{0pt plus 5em}\small\sloppy \begin{longtable}[]{@{}
  >{\hspace{0pt}\RaggedRight\arraybackslash}p{(\linewidth - 4\tabcolsep) * \real{0.3333}}
  >{\hspace{0pt}\RaggedRight\arraybackslash}p{(\linewidth - 4\tabcolsep) * \real{0.3333}}
  >{\hspace{0pt}\RaggedRight\arraybackslash}p{(\linewidth - 4\tabcolsep) * \real{0.3333}}@{}}
\toprule\noalign{}
\begin{minipage}[b]{\linewidth}\raggedright
Actividad
\end{minipage} & \begin{minipage}[b]{\linewidth}\raggedright
Canal
\end{minipage} & \begin{minipage}[b]{\linewidth}\raggedright
Frecuencia
\end{minipage} \\
\midrule\noalign{}
\endhead
\bottomrule\noalign{}
\endlastfoot
Guía de Seguridad Móvil para Estudiantes (descargable) & Portal
institucional & Publicación única + actualización anual \\
Infografías periódicas en redes sociales & Instagram / Facebook / TikTok
UNCP & Mensual \\
Cápsulas formativas en Moodle & Campus Virtual & Bimestral \\
Participación en la Semana de la Ciberseguridad & Presencial / Streaming
& Anual (agosto) \\
Alertas de amenazas activas (phishing, fraudes) & Correo institucional
estudiantil & Según evento \\
\end{longtable}\normalsize}

Los contenidos para estudiantes se enfocan en: protección de
credenciales, identificación de phishing, uso de redes Wi-Fi seguras, y
protección de dispositivos personales.

\begin{center}\rule{0.5\linewidth}{0.5pt}\end{center}

\subsection{Evaluación del Plan}\label{evaluaciuxf3n-del-plan}

\subsubsection{5.1 Indicadores de
Gestión}\label{indicadores-de-gestiuxf3n}

{\setlength{\RaggedRightRightskip}{0pt plus 5em}\small\sloppy \begin{longtable}[]{@{}
  >{\hspace{0pt}\RaggedRight\arraybackslash}p{(\linewidth - 8\tabcolsep) * \real{0.2000}}
  >{\hspace{0pt}\RaggedRight\arraybackslash}p{(\linewidth - 8\tabcolsep) * \real{0.2000}}
  >{\hspace{0pt}\RaggedRight\arraybackslash}p{(\linewidth - 8\tabcolsep) * \real{0.2000}}
  >{\hspace{0pt}\RaggedRight\arraybackslash}p{(\linewidth - 8\tabcolsep) * \real{0.2000}}
  >{\hspace{0pt}\RaggedRight\arraybackslash}p{(\linewidth - 8\tabcolsep) * \real{0.2000}}@{}}
\toprule\noalign{}
\begin{minipage}[b]{\linewidth}\raggedright
Indicador
\end{minipage} & \begin{minipage}[b]{\linewidth}\raggedright
Fórmula
\end{minipage} & \begin{minipage}[b]{\linewidth}\raggedright
Meta
\end{minipage} & \begin{minipage}[b]{\linewidth}\raggedright
Frecuencia
\end{minipage} & \begin{minipage}[b]{\linewidth}\raggedright
Fuente
\end{minipage} \\
\midrule\noalign{}
\endhead
\bottomrule\noalign{}
\endlastfoot
\textbf{Cobertura} & (N° trabajadores que completaron curso obligatorio
/ N° total de trabajadores) x 100 & \textbf{100\%} & Trimestral & Moodle
/ F-SGSI-02 \\
\textbf{Eficacia contra phishing} & (N° clics en simulacro / N° correos
enviados) x 100 & \textbf{\textless{} 10\%} & Semestral & Plataforma de
simulación \\
\textbf{Satisfacción} & Promedio de calificación de cursos (encuesta
1-5) & \textbf{\textgreater= 4.0} & Por curso & Encuesta en Moodle \\
\textbf{Competencia técnica} & N° de certificaciones de seguridad
obtenidas & \textbf{3} & Anual & Registro de certificaciones \\
\textbf{Participación en Semana de Ciberseguridad} & N° de asistentes al
evento anual & \textbf{\textgreater{} 200} & Anual & Registro de
asistencia \\
\end{longtable}\normalsize}

\subsubsection{5.2 Retroalimentación y
Mejora}\label{retroalimentaciuxf3n-y-mejora}

\begin{itemize}
\tightlist
\item
  Al finalizar cada actividad de capacitación se aplicará una encuesta
  de satisfacción.
\item
  Los resultados del simulacro de phishing se analizarán para
  identificar las áreas o perfiles más vulnerables y ajustar las
  campañas de concientización.
\item
  El Oficial de Seguridad presentará un informe semestral al Comité de
  Gobierno Digital con los resultados del plan.
\end{itemize}

\begin{center}\rule{0.5\linewidth}{0.5pt}\end{center}

\subsection{Presupuesto}\label{presupuesto}

{\setlength{\RaggedRightRightskip}{0pt plus 5em}\small\sloppy \begin{longtable}[]{@{}
  >{\hspace{0pt}\RaggedRight\arraybackslash}p{(\linewidth - 2\tabcolsep) * \real{0.5000}}
  >{\hspace{0pt}\RaggedRight\arraybackslash}p{(\linewidth - 2\tabcolsep) * \real{0.5000}}@{}}
\toprule\noalign{}
\begin{minipage}[b]{\linewidth}\raggedright
Concepto
\end{minipage} & \begin{minipage}[b]{\linewidth}\raggedright
Monto (S/)
\end{minipage} \\
\midrule\noalign{}
\endhead
\bottomrule\noalign{}
\endlastfoot
Capacitación técnica (cursos + certificaciones) & 33,000 \\
Plataforma de simulación de phishing (suscripción anual) & 5,000 \\
Materiales didácticos, diseño de infografías y videos & 4,000 \\
Semana de la Ciberseguridad (logística, ponentes, materiales) & 3,000 \\
\textbf{Total} & \textbf{45,000} \\
\end{longtable}\normalsize}

\begin{center}\rule{0.5\linewidth}{0.5pt}\end{center}

\subsection{Responsabilidades}\label{responsabilidades}

{\setlength{\RaggedRightRightskip}{0pt plus 5em}\small\sloppy \begin{longtable}[]{@{}
  >{\hspace{0pt}\RaggedRight\arraybackslash}p{(\linewidth - 2\tabcolsep) * \real{0.5000}}
  >{\hspace{0pt}\RaggedRight\arraybackslash}p{(\linewidth - 2\tabcolsep) * \real{0.5000}}@{}}
\toprule\noalign{}
\begin{minipage}[b]{\linewidth}\raggedright
Rol
\end{minipage} & \begin{minipage}[b]{\linewidth}\raggedright
Responsabilidad
\end{minipage} \\
\midrule\noalign{}
\endhead
\bottomrule\noalign{}
\endlastfoot
\textbf{Oficial de Seguridad} & Diseñar y coordinar el plan; elaborar
contenidos técnicos; analizar resultados. \\
\textbf{Unidad de RRHH} & Gestionar la inscripción del personal;
registrar asistencia (F-SGSI-02); incluir la capacitación en el plan
anual de RRHH. \\
\textbf{OTI} & Administrar la plataforma Moodle para cursos virtuales;
ejecutar simulacros de phishing. \\
\textbf{Jefes de Área} & Asegurar que su personal complete el curso
obligatorio; liberar tiempo para capacitación técnica. \\
\textbf{Imagen Institucional} & Apoyar en la difusión de campañas de
concientización en redes sociales y portal web. \\
\end{longtable}\normalsize}

\begin{center}\rule{0.5\linewidth}{0.5pt}\end{center}

\subsection{Documentos Relacionados}\label{documentos-relacionados}

{\setlength{\RaggedRightRightskip}{0pt plus 5em}\small\sloppy \begin{longtable}[]{@{}
  >{\hspace{0pt}\RaggedRight\arraybackslash}p{(\linewidth - 2\tabcolsep) * \real{0.5000}}
  >{\hspace{0pt}\RaggedRight\arraybackslash}p{(\linewidth - 2\tabcolsep) * \real{0.5000}}@{}}
\toprule\noalign{}
\begin{minipage}[b]{\linewidth}\raggedright
Código
\end{minipage} & \begin{minipage}[b]{\linewidth}\raggedright
Nombre
\end{minipage} \\
\midrule\noalign{}
\endhead
\bottomrule\noalign{}
\endlastfoot
POL-SGSI-02 & Política de Uso Aceptable, Escritorio y Teletrabajo \\
POL-SGSI-05 & Política de Seguridad para Dispositivos Móviles del
Personal \\
F-SGSI-02 & Registro de Asistencia a Capacitación \\
R-SGSI-03 & Cuadro de Mando de Seguridad (KPIs) \\
\end{longtable}\normalsize}

\begin{center}\rule{0.5\linewidth}{0.5pt}\end{center}

\textbf{Aprobado por:} Comité de Gobierno y Transformación Digital ---
UNCP

\end{document}
